%=============================================================================!
% 15.9  A PROTOCOL AND A CRITERION                                            !
%=============================================================================!
\section{A protocol and a criterion}
\label{sec:cv-protocol}

The chapter's tools combine into a procedure that takes a system from a
rough start to numbers with error bars, and a criterion, with numbers,
for when a quantity counts as converged. This section follows the liquid
from atoms placed at random to a converged run at constant pressure, with
the plot to read at each stage.

\subsection*{From random positions}

Atoms placed at random in a box overlap: 256 argon atoms dropped into the
liquid's cell start with $U = \SI{7.91e10}{\electronvolt}$ and a largest
force of \SI{2.95e12}{\electronvolt\per\angstrom}, and the first step of a
run of \SI{10}{\femto\second} would carry the atom that feels it some
\SI{4e10}{\angstrom}, across more than $10^9$ copies of the periodic box. The overlaps are
removed first by walking downhill on the energy. Moving every atom a small
distance along its own force, $\dd\vb{r}_i = h\vb{F}_i$ with $h > 0$, changes
the energy by
\begin{equation*}
   \dd U = \sum_i\grad_iU\cdot\dd\vb{r}_i = -\sum_i\vb{F}_i\cdot h\vb{F}_i
   = -h\sum_i\lvert\vb{F}_i\rvert^2 ,
\end{equation*}
which is negative unless every force vanishes: the forces point downhill,
and a step along them lowers the energy if it is short enough for the
forces not to change much along it. This is steepest descent. The function
\texttt{minimise} of \texttt{mdlab.md} sets $h$ by the atom with the
largest force, which moves the distance \texttt{step}, so that $h =
\texttt{step}/\max_i\lvert\vb{F}_i\rvert$; it keeps a trial that lowers
the energy and lengthens the next by a fifth, and refuses one that does
not and halves the distance:

\begin{code}
def minimise(
    model: EnergyForces,
    positions: ArrayLike,
    step: float = 0.1,
    force_tolerance: float = 1e-2,
    max_steps: int = 10000,
) -> tuple[NDArray, NDArray, NDArray]:
    r = np.array(positions, dtype=float)
    u, f = model(r)
    energies, largest = [u], [np.max(np.linalg.norm(f, axis=-1))]
    for _ in range(max_steps):
        if largest[-1] < force_tolerance:
            break
        trial = r + step * f / largest[-1]
        u_trial, f_trial = model(trial)
        if u_trial < u:
            r, u, f = trial, u_trial, f_trial
            energies.append(u)
            largest.append(np.max(np.linalg.norm(f, axis=-1)))
            step *= 1.2
        else:
            step *= 0.5
    return r, np.array(energies), np.array(largest)
\end{code}

\noindent From a lattice disturbed by \SI{0.2}{\angstrom} it reaches the same
minimum as ASE's FIRE optimiser, to \SI{1e-3}{\angstrom} in every position
(\texttt{test\_stats.py}). From the random start it takes 75 accepted moves
to bring the largest force below \SI{0.05}{\electronvolt\per\angstrom}, at
$U = \SI{-15.357}{\electronvolt}$ (\cref{fig:cv-protocol}(a)). The aim is
only to remove the overlaps, and any minimum close to the start will
do.

\subsection*{Temperature, then pressure}

The stages that follow each settle one thing, and each has its plot.
Velocities drawn at \SI{135}{\kelvin} for the minimised atoms give a
temperature that falls within the first picosecond to \SI{70.1}{\kelvin},
about half, as the kinetic energy shares itself with the potential energy
of atoms that sat at the bottom of their wells (\cref{sec:sm-equipartition});
CSVR brings its average over the preceding \SI{1}{\pico\second} back above
\SI{130}{\kelvin}, within \SI{5}{\kelvin} of the target, by
\SI{3.8}{\pico\second} (\cref{fig:cv-protocol}(b)), and over the last
\SI{25}{\pico\second} of \SI{50}{\pico\second} the mean is $(136.6 \pm
1.2)$\,\si{\kelvin}, 1.4 errors from the target. The running mean from the start, which lags, shows
why such a stage is discarded. Stochastic cell rescaling at
\SI{0.1}{\giga\pascal} then lets the volume move from that of
$0.8\sigma^{-3}$, \SI{12577}{\angstrom\cubed}, to the one the liquid takes
at that pressure (\cref{fig:cv-protocol}(c)), and the run continues for \SI{1}{\nano
\second}. Each quantity then goes through the criterion.

\begin{definition}
A quantity $A$, recorded at $n$ equal intervals $\dt$ over a run, is
converged to the error $\sigma^\ast$ when
\begin{enumerate}
\item the start $t_0$ that maximises $n_{\mathrm{eff}}(t_0)$ lies in the
   first half of the run (\cref{sec:cv-equilibration});
\item the $n'$ samples after it, with their own $g$ and $s$, hold at least
   100 independent samples, $n_{\mathrm{eff}} = n'/g \ge 100$
   (\cref{sec:cv-blocks});
\item the error from blocks of $5g$ samples agrees with $\sqrt{g s^2/n'}$
   within twice its own error (\cref{sec:cv-blocks});
\item the slope of the samples after $t_0$ is within two of its errors of
   zero (\cref{sec:cv-drift});
\item and $\sqrt{g s^2/n'} \le \sigma^\ast$, or else \cref{eq:cv-length}
   gives the run that would make it so (\cref{sec:cv-observables}).
\end{enumerate}
The result is reported as $\bar A \pm \sqrt{g s^2/n'}$, one standard
error, with the number of independent samples.
\end{definition}

The first check catches a run too short to locate its own start; the
second leaves at least twenty blocks of $5g$, so that the block error the
third check compares with is itself known to about a sixth,
$1/\sqrt{2\times19}$, and keeps the bias that the search for a start
brings small, as below; the third catches a slow memory that the estimate
of $g$ cut off; the fourth catches a drift or a transient left in. The
function \texttt{convergence\_report} of \texttt{mdlab.analysis.stats}
applies the first four, and the fifth when given a target, and returns
the evidence for each. The criterion says nothing about whether the run samples the
right ensemble, for which \cref{sec:cv-ensemble} gives a test, or about
the size of the box, which \cref{sec:cv-size} shows can bias a converged
number.

After \SI{300}{\pico\second} at constant pressure the report finds the
volume and the potential energy not yet converged: their starts are at
16.5 and \SI{12.0}{\pico\second}, but they hold only 64.6 and 47.8
independent samples, the potential energy now having $\tau_{\mathrm{int}}
= \SI{3.0}{\pico\second}$ over these \SI{300}{\pico\second}, as it
follows the slow swings of the volume. A
hundred independent samples need 439 and \SI{603}{\pico\second} after the
start. The pressure and the temperature pass, with 813 and 139. After
\SI{1}{\nano\second} all four pass: $V = (12310 \pm 10)$\,\si{\angstrom
\cubed} with 316 independent samples, the volume of
\cref{sec:pr-npt} within its error; $P = (0.0997 \pm 0.0004)$\,\si{\giga
\pascal} against the \SI{0.1}{\giga\pascal} asked for; $T = (135.27 \pm
0.28)$\,\si{\kelvin}; and $U = (-13.049 \pm 0.014)$\,\si{\electronvolt}.

\begin{figure}[htb]
\centering
\includegraphics{ch15_convergence/protocol}
\caption{A protocol, stage by stage, for 256 argon atoms placed at random.
(a) Steepest descent: the largest force after each accepted move, against
the tolerance (dotted). (b) \SI{50}{\pico\second} of CSVR at fixed volume:
the temperature (grey) and its running mean (teal), against
\SI{135}{\kelvin} (dotted). (c) \SI{1}{\nano\second} of stochastic cell
rescaling at \SI{0.1}{\giga\pascal}: the volume, the start found (dotted)
and the mean after it (dashed). (d) The same for the potential energy per
atom.}
\label{fig:cv-protocol}
\end{figure}

\subsection*{Does the criterion keep its promise?}

An error bar of one standard error should hold the true mean in 68\% of
runs. On series of known mean from \cref{eq:cv-ou}, 300 of each length,
the error bar the report gives, whether or not the series passes, holds it
in 0.55 of the series with 25
independent samples, 0.57 with 50, 0.68 with 100 and 0.70 with 200; the
error bar $\sqrt{g s^2/n}$ of the whole series, without the search for a
start, holds it in 0.64, 0.64, 0.71 and 0.72. The search favours a start
after which $g$ happens to come out small (\cref{sec:cv-equilibration}),
which makes the error bar small too, and only with about a hundred
independent samples does that bias fade, which sets the threshold of the
second check. With the offset of \cref{fig:cv-equilibration}(a) added to
series of \SI{200}{\pico\second}, 182 of 200 pass the criterion, and of
those 0.70 hold the true mean within one error and 0.95 within two. On the
liquid itself the evidence is thinner: of the sixty pieces of
\SI{100}{\pico\second} of $U$, $T$ and $P$ in the \SI{2}{\nano\second} run,
ten pass, and of those four hold the mean of the whole run within one
error and seven within two, fewer than 68\% and 95\% would give, which
fits pieces that pass only where $g$ came out small.

\begin{keyidea}
A protocol removes overlaps by steepest descent, brings the temperature,
then the pressure, to their targets, and discards each stage. A quantity
is converged when its start lies in the first half of the run, at least
100 independent samples follow it, blocks confirm the error, no drift
remains and the error meets its target; \texttt{convergence\_report} checks
all five. So reported, one error bar holds the true mean of test series in
about 68\% of runs; on the liquid's pieces that only just pass, less
often.
\end{keyidea}

\notebook{15}{run the criterion on any series, from mdlab or from ASE}

\begin{selfcheck}
A run of \SI{500}{\pico\second} has $g = 300$ for a quantity sampled every
\SI{10}{\femto\second}. Does it pass the second check, and if not, how long
must it be?
\end{selfcheck}
